\documentclass[11pt,a4paper]{article}
\usepackage[T1]{fontenc}
\usepackage[utf8]{inputenc}
\usepackage[ngerman]{babel}
\usepackage{lmodern}
\usepackage[margin=24mm]{geometry}
\usepackage{booktabs,tabularx,xcolor,listings,fancyhdr}
\usepackage[hidelinks]{hyperref}
\definecolor{accent}{HTML}{194B68}
\definecolor{codebg}{HTML}{F2F5F7}
\lstset{basicstyle=\ttfamily\small,backgroundcolor=\color{codebg},
  frame=none,breaklines=true,columns=fullflexible,keepspaces=true,
  aboveskip=10pt,belowskip=10pt}
\pagestyle{fancy}
\fancyhf{}
\fancyhead[L]{\small REIST / Compilervergleich}
\fancyhead[R]{\small 14. November 2026}
\fancyfoot[R]{\thepage}
\setlength{\headheight}{14pt}
\setlength{\parindent}{0pt}
\setlength{\parskip}{7pt}
\begin{document}
{\color{accent}\LARGE\bfseries REIST im Disassembler}\par
{\large Optimierungsstufen, negative Reste und vollständige Division}

\section*{1. Die Vergleichsmatrix ausführen}
Die Beispiele sind kleine Übersetzungseinheiten mit jeweils drei exportierten
C-Symbolen. Der reine Beispielbau erzeugt Assembler, Objektdateien und Listings,
ohne ein Executable zu linken. Die Matrix prüft zusätzlich das Verhalten der
genau gleichen Objektdateien mit einem temporären Testtreiber.

\begin{lstlisting}
python scripts/run_disassembly_matrix.py
\end{lstlisting}
Voraussetzungen: Python ab 3.10, GCC oder Clang und ein passendes GNU/LLVM
\texttt{objdump}. Der Matrix-Compiler muss Programme für den aktuellen Rechner
erzeugen. Für Clang beispielsweise:
\begin{lstlisting}
python scripts/run_disassembly_matrix.py --compiler clang++
    --output build/disassembly-clang
\end{lstlisting}
Den umgebrochenen Befehl als eine Zeile eingeben. Jedes Ausgabeverzeichnis muss
neu oder leer sein. Der Einstieg ist danach
\texttt{build/disassembly-matrix/REPORT.md} beziehungsweise das entsprechende
Clang-Verzeichnis.

\begin{tabularx}{\linewidth}{@{}lX@{}}
\toprule
Profil & Optionen und Zweck \\
\midrule
O0 & Ohne Optimierung; Aufrufe nicht eingebetteter Helfer nachvollziehen. \\
O1, O2, O3 & Zunehmende Optimierung, Auto-Vektorisierung ausgeschaltet. \\
Os & Auf Codegröße optimieren, Auto-Vektorisierung ausgeschaltet. \\
O3\_auto & O3 mit erlaubter Auto-Vektorisierung; SIMD ist nicht garantiert. \\
\bottomrule
\end{tabularx}

Alle Profile deaktivieren Loop-Unrolling und verwenden kein LTO. Ohne
\texttt{--native} gilt das ISA-Standardziel des Compilers; mit dieser Option
gilt \texttt{-march=native} für alle Profile. Compiler-Version, Target, Optionen,
Befehle und Datei-Hashes stehen in den Manifesten.

\textbf{Dateien zum Vergleichen:}
Pro Beispiel entstehen \texttt{.s}, \texttt{.o} und \texttt{.disasm.txt}.
Unter \path{<Profil>/functions/} liegt zusätzlich ein Listing pro Funktion.
\texttt{SUMMARY.csv} enthält Instruktions-, DIV-, Aufruf- und Sprungzahlen.
\texttt{remainder\_samples.csv} enthält geprüfte Quotienten und Reste.
Die CI veröffentlicht vollständige GCC- und Clang-Matrizen als Artefakte.

\textbf{Nur ein Objekt bauen:}
\begin{lstlisting}
python scripts/build_disassembly_examples.py
    --example 09_division_257 --optimization Os
    --output build/division-257-Os
\end{lstlisting}

\newpage
\section*{2. Welche Arithmetik vergleichen wir?}
Die Beispiele 01 bis 05 verwenden vorbereitete 32-Bit-Reste für modulare
Addition. Verglichen werden \texttt{\%}, eine beschränkte nicht-negative
Korrektur und die zentrierte Korrektur. Es gibt Laufzeit- und Konstantenmoduli,
eine Zweierpotenz als Kontrolle sowie serielle und unabhängige Array-Updates.
Diese Additionsfunktionen übernehmen keine Eingabevalidierung.

Die Beispiele 06 bis 09 behandeln beliebige 64-Bit-Dividenden. Sie vergleichen
die drei Restkonventionen sowohl beim reinen Rest als auch bei Quotient und
Rest zusammen. Die REIST-Varianten verwenden die öffentliche Implementierung
aus \texttt{include/reist\_mod.hpp}.

\begin{tabularx}{\linewidth}{@{}lX@{}}
\toprule
Datei & Inhalt \\
\midrule
06\_signed\_remainders\_runtime.cpp & Drei Reste; B als Parameter. \\
07\_signed\_remainders\_257.cpp & Drei Reste; B konstant 257. \\
08\_division\_runtime.cpp & Drei Quotient-/Restpaare; B als Parameter. \\
09\_division\_257.cpp & Drei Quotient-/Restpaare; B konstant 257. \\
\bottomrule
\end{tabularx}

Bei den allgemeinen Laufzeitvarianten gilt dieselbe Validierung $B>0$.
Andernfalls werfen alle Varianten \texttt{std::invalid\_argument}, bevor
gerechnet oder ein Ergebnis geschrieben wird. Exception- und
Unwind-Unterstützung bleiben aktiviert. Die vollständigen Divisionen schreiben
Quotient und Rest über zwei verschiedene gültige Ausgabezeiger.

Für jede Darstellung gilt $T=qB+r$, aber die Restintervalle unterscheiden sich:
C++ dividiert gegen null; ein Rest hat das Vorzeichen des Dividenden oder ist
null. Euklidische Reste liegen in $[0,B)$. REIST verwendet ganzzahlig
$[-\lfloor B/2\rfloor,\lceil B/2\rceil)$.

\begin{center}
\begin{tabular}{@{}rrccc@{}}
\toprule
T & B & C++ $(q,r)$ & Euklidisch $(q,r)$ & REIST $(q,r)$ \\
\midrule
-260 & 257 & $(-1,-3)$ & $(-2,254)$ & $(-1,-3)$ \\
-129 & 257 & $(0,-129)$ & $(-1,128)$ & $(-1,128)$ \\
129 & 257 & $(0,129)$ & $(0,129)$ & $(1,-128)$ \\
128 & 256 & $(0,128)$ & $(0,128)$ & $(1,-128)$ \\
\bottomrule
\end{tabular}
\end{center}

\textbf{Korrektheit:} Pro Profil werden alle 27 Funktionen geprüft. Für die
allgemeine Division sind es 548\,819 Eingabe-/Modulpaarungen: kleine Werte
vollständig, reproduzierbare Zufallswerte, gerade und ungerade Moduli sowie
die Grenzen \texttt{INT64\_MIN} und \texttt{INT64\_MAX}. Hinzu kommen ungültige
Moduli sowie Additions- und Schleifentests. Die 170 CSV-Beispielzeilen müssen
über alle Profile identisch sein. Der temporäre Testtreiber wird entfernt;
die kleinen Vergleichsobjekte bleiben erhalten.

Die allgemeine REIST-Division reduziert beliebige Eingaben. Die schnelle
Addition mit bereits vorbereitetem zentriertem Zustand ist ein eigener Fall.
Eine Aussage über einen dieser Fälle darf nicht auf den anderen übertragen werden.

\newpage
\section*{3. Geprüfte Beobachtungen mit GCC 15.2}
Lokaler Lauf vom 14.11.2026 mit GCC 15.2 für
\texttt{x86\_64-w64-mingw32}, ohne \texttt{-march=native}:
Alle sechs Profile bestanden die Tests. Die neun O2-Objekte sind 961 bis
3681 Bytes groß; über die ganze Matrix sind es 945 bis 6112 Bytes.
Diese Größen enthalten COFF-Metadaten und gegebenenfalls Fehlerbehandlung.

Bei konstantem B=257 erzeugt GCC für \texttt{reist\_division\_constant} unter
O2 Multiplikation und Shifts zur Quotientenbildung. Ein Auszug aus dem
vollständigen Listing:
\begin{lstlisting}
movabs rax,0x7f807f807f807f81
mov    r9,rdx
imul   rcx
mov    rax,rcx
sar    rax,0x3f
sar    rdx,0x7
sub    rdx,rax
\end{lstlisting}
Danach folgen Restbildung und die Korrekturen von Quotient und Rest.
Unter Os verwendet dieselbe Funktion dagegen eine Hardwaredivision:
\begin{lstlisting}
mov    r9d,0x101
mov    rax,rcx
mov    rcx,rdx
cqo
idiv   r9
\end{lstlisting}
Die Laufzeitvariante benötigt unter O2 ein \texttt{idiv} für Quotient und Rest
zusammen. Bei Os kann zusätzlich die Division durch zwei zur Berechnung der
REIST-Grenzen als \texttt{idiv} erscheinen. Auch beim beschränkten
Additionsbeispiel kann diese Grenzberechnung eine DIV-Instruktion erzeugen.
Die Operanden zeigen, ob durch B oder durch zwei dividiert wird.

\textbf{Listings richtig lesen:}
Die Zählung erfasst den benannten Funktionskörper einschließlich Alignment,
jedoch keine aufgerufenen Helfer, Tailcall-Ziele oder ausgelagerten Fehlerpfade.
Besonders bei O0 kann die eigentliche Division in einem Helfer liegen.
Die vollständigen Objektlistings mit Relokationen erlauben, diesen Aufrufen
zu folgen. Die Klassifikation der Mnemonics unterstützt x86 und AArch64.

Weniger Instruktionen bedeuten nicht automatisch geringere Laufzeit. Ein
fehlendes lokales DIV beweist weder divisionsfreie Gesamtausführung noch
konstantes Zeitverhalten. Andere Compiler, Versionen und Zielarchitekturen
können anders entscheiden; genau dafür wird die Matrix neu erzeugt.

\small Quellen im Repository:\par
\path{examples/disassembly/README.md}\\
\path{scripts/run_disassembly_matrix.py}\\
\path{tests/test_disassembly_examples.cpp}\par
Die Anleitung reproduzieren: Nach Anlegen von \path{build/pdf/} den folgenden
Befehl ausführen und das PDF nach \path{docs/papers/} kopieren.
\begin{lstlisting}[basicstyle=\ttfamily\footnotesize]
pdflatex -output-directory=build/pdf docs/papers/REIST-Compiler-Vergleich.tex
\end{lstlisting}
\end{document}
